\begin{lemma}[The full-tree route retains its connecting seed]
Fix a hierarchy counterexample $C$ on $T$, its minimum context, and
colorings $\varphi,\sigma$. Suppose there are
\noderef{P2752_T2_PrefixData} $F$ on a penultimate exact prefix $B^+$
for $\sigma$, every displayed edge of $T$ has the same value under both
colorings, and $\sigma$ is $(T_n,D_n,\varphi)$-stable on the old terminal
ladder. Then there exists a corresponding ETT certificate on the full
exact $T$ under $\sigma$, and its replacement ETT has maximum property.
The full-tree certificate is not required to equal the prefix certificate
in $F$.
\end{lemma}
\begin{proof}
Write $n$ for the old rung and $T_{n+1}$ for the final certified-series
tree. The above-penultimate, strictness and below-final fields of
\noderef{ExtendedTashkinovTree} give
$T_n\prec T\preccurlyeq T_{n+1}$ and identify the generating coloring
with the terminal series coloring. Inspect the actual certificate
$\mathit{ett.series.certified}(n)$ from \noderef{TashkinovSeries}.
It has the constructors of \noderef{TashkinovSeriesStep}. Its initial
constructor is impossible, since its index is zero whereas $n\ge1$.

For a PE certificate, its closure field, expanded by
\noderef{TashkinovClosure}, supplies the augmentation route
$T_n\longrightarrow T_{n+1}$ under $\varphi$. The exact relations above
allow \noderef{TashkinovAugmentationSequenceRestrict} to truncate it at
$T$. For an SE certificate retain its actual choice edge $f$ and its
actual adjoined tree $S$. For an RE certificate retain instead its actual
revisiting edge $f$ and its actual adjoined tree $S$. In either case the
constructor gives $F_n=\{f\}$, $\noderef{TreeSequenceAdjoinEdge}(T_n,f,S)$,
and a closure route $S\longrightarrow T_{n+1}$ under $\varphi$.
The adjoin relation gives $|S|=|T_n|+1\le|T|$.
By \noderef{TashkinovAugmentationSequenceInitialSegment}, the closure
route places $S$ as an exact prefix of $T_{n+1}$. Both $S$ and $T$ have
the same entries as $T_{n+1}$ through the indices of $S$, so
\noderef{TreeSequenceInitialSegment} gives $S\preccurlyeq T$.
Restricting this closure route thus gives $S\longrightarrow T$.
In particular we have retained the actual exact singleton and adjoined
seed, not an arbitrary one-edge extension of $T_n$.

Transport either restricted route by
\noderef{TashkinovAugmentationSequenceColoringTransport}, with $B^+$ as
penultimate prefix and $T$ as target. Its prefix and unit-length premises
are fields of $F$. Its fixed-edge premise holds at every required index
by the assumed displayed-edge equality on the full $T$. Its missing-set
premise is the vertex clause of the prefix invariant in $F$, using
\noderef{TreeColoringInvariant}. Thus in PE the identical route starts
at $T_n$, and in SE or RE it starts at the identical $S=T_n+f$ with the
identical connecting singleton. The last displayed edge of $T$ is included
in this transport, because fixed-edge equality was assumed on all of $T$.

Apply \noderef{Td} to the original ETT, its positive rung and maximum
property from \noderef{HierarchyCounterexample}, and the entire
six-assertion lower-rung field of
\noderef{HierarchyMinimumCounterexampleContext}. Properness and exact
one-uncoloredness of $\sigma$ are fields of $F$, and terminal-ladder
stability is the stated hypothesis, rewritten using the generating-color
equation. This yields \noderef{ModColoring}(\mathit{ett},\sigma).
Its universal clause is applicable to $T$ with $T_n\prec T$. In PE its
first eligibility alternative is supplied by the transported PE route;
in SE or RE its second alternative is supplied by the terminal kind,
actual singleton $\{f\}$, adjoin relation and transported seed route.
It therefore supplies an actual full-tree
\noderef{CorrespondingExtendedTashkinovTree}. Apply
\noderef{MaximumPropertyTransport} to this mod-coloring certificate,
this full-tree corresponding certificate, and the original maximum
property. It gives maximum property of the replacement ETT. Nothing
in this argument identifies this full certificate with the independently
stored prefix certificate.
\end{proof}
